\documentclass[11pt,a4paper]{article}
\usepackage[utf8]{inputenc}
\usepackage[T1]{fontenc}
\usepackage{tgpagella}
\usepackage[spanish]{babel}
\usepackage{geometry}
\usepackage{booktabs}
\usepackage{xcolor}
\usepackage{titlesec}
\usepackage{fancyhdr}
\usepackage[parfill]{parskip}
\usepackage[expansion=false]{microtype}
\usepackage{hyperref}
\usepackage{setspace}
\usepackage{needspace}
\usepackage{etoolbox}
\usepackage{graphicx}

\widowpenalty=10000
\clubpenalty=10000
\displaywidowpenalty=10000
\raggedbottom
\sloppy
\hbadness=10000
\hfuzz=2pt

\geometry{top=3.0cm,bottom=3.0cm,left=3.5cm,right=3.5cm}
\setlength{\parskip}{0.65em}
\setlength{\parindent}{0em}
\setlength{\headheight}{14pt}

\definecolor{azulai}{RGB}{30,80,140}
\definecolor{doradoai}{RGB}{140,90,0}
\definecolor{grisai}{RGB}{70,70,70}

\titleformat{\section}{\Large\bfseries\color{azulai}}{}{0em}{}[{\color{azulai}\titlerule[0.5pt]}]
\titlespacing*{\section}{0pt}{1.8em}{0.8em}

\titleformat{\subsection}{\large\bfseries\color{doradoai}}{}{0em}{}
\titlespacing*{\subsection}{0pt}{1.4em}{0.5em}

\titleformat{\subsubsection}{\normalsize\bfseries\color{grisai}}{}{0em}{}
\titlespacing*{\subsubsection}{0pt}{1.0em}{0.3em}

% Evita títulos huérfanos al final de página
\pretocmd{\section}{\Needspace{12\baselineskip}}{}{}
\pretocmd{\subsection}{\Needspace{9\baselineskip}}{}{}
\pretocmd{\subsubsection}{\Needspace{6\baselineskip}}{}{}

\pagestyle{fancy}
\fancyhf{}
\rhead{\textcolor{grisai}{\small Maite Barragan — Arquitectura Interna}}
\lhead{\textcolor{grisai}{\small Luna · Casa 4 · Casa 6}}
\cfoot{\textcolor{grisai}{\small\thepage}}
\renewcommand{\headrulewidth}{0.3pt}

\hypersetup{colorlinks=true,linkcolor=azulai,urlcolor=azulai}
\setstretch{1.45}
\tolerance=2500
\emergencystretch=8em

\begin{document}

% ── Portada ──────────────────────────────────────────────────────────────────
\begin{titlepage}
  \centering
  \vspace*{1.5cm}
  {\Huge\bfseries\color{azulai} Luna · Casa 4 · Casa 6}\\[0.5cm]
  {\large\color{grisai} Arquitectura Interna}\\[0.3cm]
  {\small\itshape\color{grisai} Un método para sostener cuerpo, energía y vida con coherencia}\\[2cm]
  {\huge\color{doradoai} Maite Barragan}\\[1.5cm]
  {\Large 07/01/1971 \quad 12:15}\\[0.3cm]
  {\Large Pamplona, España}\\[0.3cm]
  {\normalsize Lat: 42.8157° \quad Lon: -1.6522° \quad UTC+01:00}\\[0.3cm]
  {\normalsize Ascendente: Aries 6°18' \quad
    MC: Capricornio 3°11'}\\[2cm]

  \begin{tabular}{lll}
    \textbf{Luna:} & Tauro & Casa 2 \\
    \textbf{Casa 4:} & Cáncer en cúspide & Regente: Luna \\
    \textbf{Casa 6:} & Leo en cúspide & Regente: Sol \\
  \end{tabular}\\[2cm]

  \vfill
  {\small Generado el 20/05/2026}
\end{titlepage}

\tableofcontents
\clearpage

% ── Datos de referencia ───────────────────────────────────────────────────────
\section{Datos de referencia}

\begin{center}
\begin{tabular}{llll}
  \toprule
  \textbf{Eje} & \textbf{Signo} & \textbf{Posición} & \textbf{Datos} \\
  \midrule
  Luna & Tauro & Casa 2 & 27°50' \\
  Casa 4 (cúspide) & Cáncer & --- & Regente: Luna \\
  Casa 6 (cúspide) & Leo & --- & Regente: Sol \\
  Ascendente & Aries & --- & 6°18' \\
  Medio Cielo & Capricornio & --- & 3°11' \\
  \bottomrule
\end{tabular}
\end{center}

\vspace{0.4cm}
\textbf{Planetas en Casa 4:} ninguno\\
\textbf{Planetas en Casa 6:} Plutón, Lilith

\vspace{0.5cm}
\textbf{Aspectos de la Luna:}

\begin{center}
\begin{tabular}{lllll}
  \toprule
  & \textbf{Asp.} & \textbf{Planeta} & \textbf{Tipo} & \textbf{Orbe} \\
  \midrule
  Luna & ⚻ & Mercurio & Quincuncio & 0.0° \\
  Luna & ☍ & Júpiter & Oposición & 0.9° \\
  Luna & △ & Plutón & Trígono & 1.9° \\
  Luna & ☍ & Venus & Oposición & 2.5° \\
  Luna & □ & Nodo Norte & Cuadratura & 3.6° \\
  Luna & □ & Nodo Sur & Cuadratura & 3.6° \\
  Luna & ☍ & Neptuno & Oposición & 4.3° \\
  Luna & ☍ & Marte & Oposición & 7.7° \\
  \bottomrule
\end{tabular}
\end{center}

\vspace{0.8cm}
\begin{center}
\includegraphics[width=0.80\textwidth]{C:/Users/Izas/claude/astrologia/pruebas/Maite_Barragan_Luna_C4_C6_rueda.png}
\end{center}


\vspace{0.5cm}
\Needspace{5\baselineskip}

% ── Interpretación ────────────────────────────────────────────────────────────
\section{Interpretación — Arquitectura Interna}

\Needspace{8\baselineskip}
\begin{center}
{\small\itshape
No se trata de explicarte. Se trata de mostrar cómo funciona tu regulación interna,\\
dónde puede haber desregulación y qué necesitas para sostenerte.
}
\end{center}

\vspace{0.8cm}
\Needspace{5\baselineskip}

% ── 1. Luna ───────────────────────────────────────────────────────────────────
\subsection{1. Tu Luna — Función emocional}

\subsubsection*{Luna en Tauro — Casa 2}

Los procesos emocionales necesitan tiempo para asentarse. Lo que te afecta no desaparece rápido: se queda dentro mientras tu cuerpo intenta procesarlo a su ritmo.

Muchas veces notas antes la tensión física que la emoción en sí. Puede sentirse como cansancio, necesidad de descanso, rigidez o sensación de pesadez acumulada.

Los cambios bruscos suelen desorganizarte más de lo que parece desde fuera, porque necesitas continuidad para recolocarte internamente. Cuando existe demasiada presión para acelerar procesos emocionales, la respuesta suele ser todavía más resistencia.

La estabilidad, los ritmos previsibles y la sensación de seguridad concreta te ayudan a recuperar regulación. En cambio, los cambios repentinos, la presión externa o la sensación de perder estabilidad tienden a generar desorganización emocional.

Tu estabilidad emocional está muy ligada a la sensación de seguridad concreta. Cuando aparece incertidumbre económica, cambios materiales o dudas sobre tu propio valor, todo dentro suele moverse rápidamente. No se vive solo como una preocupación mental. Puede aparecer tensión física, necesidad de controlar más las cosas o sensación de perder suelo interno. Necesitas cierta continuidad y estabilidad para relajarte de verdad. Cuando todo cambia demasiado rápido, puedes permanecer en alerta aunque desde fuera parezca que sigues funcionando con normalidad. Te ayuda sentir que existe una base suficiente sobre la que apoyarte y que la vida tiene cierta estabilidad. En cambio, las pérdidas, la inseguridad o sentir que no tienes dónde sostenerte suelen generar mucho desorden interno.

Existe un ajuste constante entre sentir y entender. A veces necesitas hablar para ordenar lo que ocurre y otras veces hablar demasiado puede alejarte de lo que realmente sientes. No existe una única fórmula estable. Tu equilibrio requiere aprender cuándo poner palabras y cuándo permanecer en contacto con la emoción sin intentar explicarla enseguida.

Puede haber oscilación entre necesidad emocional y exceso. A veces buscas amplitud, sentido o afecto en una escala mayor de la que el entorno realmente puede sostener. Y cuando eso no llega, puede aparecer una sensación de vacío muy grande o de desilusión difícil de llenar. Tu bienestar depende mucho de aprender medida emocional sin reducir la amplitud natural de lo que sientes.

Tienes una capacidad muy profunda de transformación emocional. Cuando algo te afecta, puedes atravesarlo con una fuerza interna enorme aunque el proceso no siempre sea sencillo. Existe capacidad real para regenerarte desde dentro y reconstruirte después de experiencias intensas. La emoción aquí no se queda solo en impacto: puede convertirse en comprensión, fuerza y transformación profunda cuando encuentra un cauce adecuado.

Puede haber tensión entre lo que necesitas emocionalmente y la forma en que buscas afecto. A veces intentas recibir amor adaptándote, agradando o cuidando más de lo que realmente puedes sostener. Cuando eso ocurre, una parte de ti queda esperando algo que no termina de pedir directamente. La sensación de equilibrio aparece cuando el afecto deja de implicar renunciar a lo que tú también necesitas.

Tu mundo emocional puede entrar en tensión con la dirección de crecimiento que la vida intenta abrir para ti. A veces lo que resulta emocionalmente familiar o cómodo no es necesariamente lo que más ayuda a avanzar. Existe tendencia a volver a patrones conocidos porque generan sensación de seguridad, incluso cuando ya no acompañan tu evolución actual. Tu regulación mejora cuando puedes distinguir entre seguridad conocida y crecimiento real.

Hay patrones emocionales antiguos que pueden generar tensión interna en el presente. A veces reaccionas desde una memoria emocional que ya no corresponde completamente a la situación actual. Algo se activa ahora, pero la respuesta puede venir de un lugar mucho más antiguo dentro de ti. El aprendizaje está en reconocer cuándo algo pertenece realmente al presente y cuándo responde a formas anteriores de protección emocional.

\vspace{0.8cm}
\Needspace{5\baselineskip}

% ── 2. Casa 4 ─────────────────────────────────────────────────────────────────
\subsection{2. Tu Casa 4 — Raíz interna}

\subsubsection*{Cáncer en la cúspide — Regente: Luna}

Tu mundo interno responde profundamente al clima emocional del entorno cercano. Cuando existe tensión en casa, en la familia o en las personas que sientes como propias, el cuerpo suele notarlo enseguida. Necesitas sentir que existe un lugar emocionalmente seguro donde puedas bajar la guardia de verdad. Cuando aparece desbordamiento emocional, tiendes a buscar refugio en lo conocido para recuperar sensación de protección. Volver a lo familiar suele ayudarte a sentir estabilidad otra vez. Te ayuda la intimidad, la cercanía emocional y sentir que tienes un refugio real. En cambio, la inestabilidad afectiva, la tensión familiar o sentir que no tienes dónde apoyarte suele mantenerte en alerta durante demasiado tiempo.

El regente de la Casa 4 es la Luna, en Tauro y en Casa 2. La base interna tiende a organizarse desde los recursos y el valor propio, con la cualidad concreta y sostenida de Tauro.

\vspace{0.8cm}
\Needspace{5\baselineskip}

% ── 3. Casa 6 ─────────────────────────────────────────────────────────────────
\subsection{3. Tu Casa 6 — Regulación cotidiana}

\subsubsection*{Leo en la cúspide — Regente: Sol}

Tu energía cotidiana necesita algo que sientas verdaderamente propio. Cuando toda la vida se reduce a obligación, rutina o gestión, algo dentro empieza a apagarse poco a poco. Puedes seguir funcionando durante mucho tiempo, pero sin sensación real de vitalidad. Necesitas espacios donde puedas expresarte, crear o sentir conexión auténtica contigo. Cuando pasas demasiado tiempo lejos de esa parte interna, el cuerpo empieza a vaciarse antes incluso de que puedas ponerle nombre a lo que te ocurre. Te ayudan la creatividad, la expresión auténtica y sentir conexión personal con lo que haces. En cambio, la rutina mecánica, el exceso de obligación o no tener espacio propio suele apagar progresivamente la energía.

El regente de la Casa 6 es el Sol, en Capricornio y en Casa 10. La regulación cotidiana tiende a organizarse desde la vocación y el espacio público, con la cualidad concreta y sostenida de Capricornio.

Tu relación con el cuerpo y los hábitos suele vivirse con mucha intensidad. Las rutinas superficiales o automáticas rara vez consiguen sostenerte durante mucho tiempo, porque necesitas sentir transformación real y profundidad en lo que haces. Cuando atraviesas una crisis física o emocional, tu capacidad de regeneración puede ser enorme. Pero antes de regenerarte, muchas veces atraviesas procesos muy intensos que lo remueven todo profundamente. Tu cuerpo no suele funcionar a medias.

Tu cuerpo no siempre responde bien a normas impuestas desde fuera. Existe una parte muy instintiva en tu manera de relacionarte con la salud, el descanso y el cuidado cotidiano. Muchas veces necesitas descubrir personalmente qué es lo que realmente te funciona y qué no encaja contigo aunque aparentemente sea lo correcto. Cuando intentas forzarte continuamente a seguir modelos ajenos, el cuerpo acaba reaccionando de alguna manera. Escuchar lo que necesitas de verdad suele ser mucho más importante que intentar cumplir expectativas externas.

\vspace{0.8cm}
\Needspace{5\baselineskip}

% ── 4. Integración ────────────────────────────────────────────────────────────
\subsection{4. Integración — Cómo funcionan las tres capas juntas}

Tu Luna en Tauro muestra cómo reaccionas emocionalmente cuando algo te afecta.

Es el primer movimiento: lo que aparece antes de que puedas ordenarlo, explicarlo o decidir qué hacer con ello.

Después, esa emoción busca un lugar donde apoyarse. Ahí entra tu Casa 4, con Cáncer en la cúspide: la parte de ti que necesita sentir base, refugio y seguridad interna.

Si esa base está disponible, lo que sientes puede moverse con más facilidad.
Si no lo está, la emoción se queda más tiempo en estado de alerta.

Y lo que no encuentra apoyo interno termina llegando al cuerpo. Ahí entra tu Casa 6, con Leo en la cúspide: tus hábitos, tu energía diaria, tu descanso y la forma en que el cuerpo sostiene lo que vives.

Por eso, cuando algo se desordena emocionalmente, no se queda solo en lo emocional. Puede acabar apareciendo como cansancio, pérdida de ritmo, dificultad para descansar o sensación de no poder sostener el día con la misma energía.

En tu caso, las tres capas tienen elementos distintos: Tierra en la Luna, Agua en la Casa 4 y Fuego en la Casa 6.

Esto significa que tu emoción, tu base interna y tu cuerpo no siempre piden lo mismo.

Lo que te ayuda emocionalmente puede no ser lo que tu cuerpo necesita ese día. Lo que te da seguridad interna puede no coincidir con el ritmo que sostiene tu rutina.

Por eso es importante no buscar una única respuesta para todo. A veces una parte de ti necesita expresión, otra necesita refugio y otra necesita descanso, orden o movimiento.

Además, Plutón y Lilith en la Casa 6 influye en cómo tu cuerpo sostiene la vida cotidiana.

No se trata solo de hábitos o rutina. Se trata de cómo organizas tu energía, cómo respondes al cansancio y qué ocurre cuando el cuerpo empieza a pedir ajuste.

Esos planetas muestran qué tipo de fuerza, sensibilidad o tensión entra en tu regulación diaria.

Los aspectos de tensión con Mercurio, Júpiter y Venus hacen que el mundo emocional tenga menos margen en ciertos momentos.

Cuando algo te activa, puede haber menos tiempo entre sentir y reaccionar. La emoción puede llegar más intensa, más mezclada o más difícil de ordenar.

Esto no significa que sí o sí vayas a desregularte. Significa que necesitas reconocer antes las señales iniciales, porque cuando la activación ya ha subido mucho, cuesta más volver al centro.

Al mismo tiempo, los aspectos de apoyo con Plutón ofrecen recursos reales.

Cuando consigues encontrar un primer punto de estabilidad, esos aspectos ayudan a sostenerlo.

No eliminan la tensión, pero pueden facilitar que vuelvas a organizarte sin quedarte completamente dentro de lo que se activó.

Lo más difícil no suele ser que una sola parte se altere.

Lo más difícil aparece cuando las tres capas se mueven al mismo tiempo:
la Luna en Tauro, con su forma concreta de sentir;
la Casa 4 en Cáncer, con su necesidad de base y seguridad interna;
y la Casa 6 en Leo, con el cuerpo intentando sostenerlo todo en la vida diaria.

Cuando esto ocurre, puedes notar que no sabes por dónde empezar: emocionalmente hay activación, internamente falta suelo y el cuerpo empieza a mostrar cansancio, tensión o pérdida de ritmo.

El trabajo no consiste en resolverlo todo a la vez.

Consiste en encontrar la primera capa disponible: una emoción que pueda expresarse, un entorno que pueda darte más seguridad, o una acción corporal sencilla que te ayude a volver poco a poco a la estabilidad.

\vspace{0.8cm}
\Needspace{5\baselineskip}

% ── 5. Sostén y desregulación ─────────────────────────────────────────────────
\subsection{5. Sostén y desregulación}

Tu Luna en Tauro se desregula cuando lo que sientes no encuentra salida.

Puede ser porque no puedes expresarlo, porque no hay espacio para moverte, porque el entorno no lo recibe o porque intentas sostenerlo demasiado tiempo dentro.

No siempre hace falta que ocurra algo muy intenso.

A veces basta con acumular durante demasiado tiempo una emoción que no ha tenido lugar, palabra, movimiento o descanso.

Los aspectos de tensión con Mercurio, Júpiter y Venus hacen que ese margen sea más estrecho.

Cuando algo te activa, puede haber menos tiempo entre sentir y reaccionar.

La emoción puede llegar más fuerte, más mezclada o más difícil de ordenar. Por eso es importante reconocer las primeras señales antes de llegar al punto de desborde.

Tu Casa 4 con Cáncer pierde estabilidad cuando dejas de sentir base interna.

No siempre tiene que haber un conflicto claro.

A veces basta con que cambie el clima emocional del entorno, con que una relación cercana se tense, con que no puedas descansar de verdad o con que el espacio íntimo deje de sentirse seguro.

Cuando eso ocurre, no solo cambia el ánimo. Cambia la sensación de tener suelo por dentro.

Tu Casa 6 con Leo muestra la desregulación en el cuerpo y en la vida cotidiana.

Puede aparecer como cansancio, pérdida de ritmo, tensión física, sueño menos reparador o dificultad para sostener hábitos que normalmente sí puedes sostener.

Muchas veces el cuerpo avisa antes de que puedas explicar qué está pasando.

Cuando el cuerpo empieza a perder continuidad, conviene mirar también qué emoción no ha encontrado salida y qué parte de tu base interna ha dejado de sentirse segura.

Los momentos más difíciles aparecen cuando las tres capas se mueven a la vez.

La Luna en Tauro se activa.
La Casa 4 en Cáncer deja de sentirse como base.
Y la Casa 6 en Leo empieza a mostrarlo en el cuerpo, el descanso o la rutina.

Ahí puedes sentir que no sabes por dónde empezar.

No se trata de resolverlo todo a la vez.
Se trata de encontrar la primera capa disponible: expresar algo, descansar, ordenar una pequeña rutina, moverte, pedir espacio o recuperar una mínima sensación de seguridad.

Los aspectos de apoyo con Plutón funcionan como recursos cuando consigues recuperar un primer punto de estabilidad.

No eliminan la desregulación, pero pueden ayudarte a reorganizarte antes.

Cuando aparece una mínima base, esos aspectos facilitan que no te quedes dentro de lo que se activó.

\vspace{0.8cm}
\Needspace{5\baselineskip}

% ── 6. Orientación final ──────────────────────────────────────────────────────
\subsection{6. Orientación final}

\subsubsection*{Qué observar}

Empieza observando qué situaciones activan tu Luna en Tauro.

No para corregirte.
No para controlar lo que sientes.

Solo para reconocer antes cómo empieza la activación.

¿Qué ocurre justo antes de sentirte que la situación te sobrepasa?
¿Qué cambia en el entorno, en los vínculos o en el ritmo diario?

Muchas veces la desregulación no empieza en el momento del desborde. Empieza bastante antes: cuando tu Casa 4 en Cáncer deja de sentirse segura o cuando tu Casa 6 en Leo pierde continuidad y el cuerpo empieza a acumular tensión.

\subsubsection*{Qué puede desregularte}

En tu caso, la desregulación suele seguir una secuencia.

Primero, tu Luna en Tauro busca algo que necesita para estabilizarse y no lo encuentra.

Después, tu Casa 4 deja de funcionar como base interna. Puede ser por tensión emocional, por cambios en el entorno, por exceso de carga o simplemente porque llevas demasiado tiempo sosteniendo más de lo que puedes absorber.

Y finalmente, la Casa 6 en Leo empieza a mostrarlo en el cuerpo:
menos energía,
menos descanso,
menos capacidad de sostener hábitos o continuidad.

Cuando las tres capas coinciden, puedes sentir que todo cuesta más de lo habitual aunque desde fuera no parezca estar ocurriendo nada tan grave.

\subsubsection*{Qué te estructura y sostiene}

Lo que te ayuda a regularte no tiene que ser enorme.

De hecho, muchas veces lo más importante son cosas pequeñas que se sostienen en el tiempo.

Una mínima rutina corporal.
Un espacio de descanso real.
Un entorno que se sienta más seguro.
Una forma de expresar lo que ocurre antes de acumularlo demasiado.

Cuando tu Casa 4 siente suelo, tu Luna necesita menos esfuerzo para estabilizarse.

Y cuando tu Casa 6 mantiene cierta continuidad —aunque sea pequeña— el cuerpo puede absorber mejor lo que emocionalmente estás atravesando.

No se trata de no sentir.
Se trata de que lo que sientes no tenga que terminar descargándose siempre en agotamiento, tensión o pérdida de equilibrio.

\subsubsection*{Práctica}

Tu regulación necesita movimiento.

Dedica unos minutos al día a mover el cuerpo conscientemente.

No hace falta intensidad extrema. Hace falta salida.

Caminar rápido, movilizar el cuerpo, sacudir tensión o activar la respiración puede ayudar a que la energía no se quede acumulada internamente.

Cuando el cuerpo puede descargar, la emoción también encuentra más espacio para ordenarse.

\Needspace{8\baselineskip}
\vspace{0.6cm}

\begin{center}
{\small\itshape\color{grisai}
La astrología se usa aquí como lenguaje simbólico de observación, no como definición de la persona.\\
Lo que tienes delante es un mapa de posibilidades, no un destino fijo.
}
\end{center}

\end{document}
